\documentclass[12pt]{article}
\usepackage[top=1in, bottom=1in, left=.75in, right=.75in]{geometry}
\usepackage{amsmath}
\usepackage{fancyhdr}
\usepackage{graphicx}
\usepackage{txfonts}
\usepackage{multicol,coordsys}
\usepackage[scaled=0.86]{helvet}
\renewcommand{\emph}[1]{\textsf{\textbf{#1}}}
\usepackage{anyfontsize}
\usepackage[shortlabels]{enumitem}
\usepackage{tikz,pgfplots,mathrsfs}
\usetikzlibrary{calc, backgrounds}

\pgfplotsset{compat = newest, my style/.append style={axis x line=middle, axis y line=
middle, xlabel={$x$}, ylabel={$y$},axis equal}}

%yticklabels={,,} , xticklabels={,,}

% \setmainfont{Times}
% \def\sansfont{Lucida Grande Bold}
\parindent 0pt
\parskip 4pt
\pagestyle{fancy}
\fancyfoot[C]{\emph{\thepage}}
\fancyhead[L]{\ifnum \value{page} > 1\relax\emph{Math 252: Midterm Exam 1}\fi}
\fancyhead[R]{\ifnum \value{page} > 1\relax\emph{Spring 2026}\fi}
\headheight 15pt
\renewcommand{\headrulewidth}{0pt}
\renewcommand{\footrulewidth}{0pt}
\let\ds\displaystyle
\def\continued{{\emph {Continued....}}}
\def\continuing{{\emph {Problem \arabic{probcount} continued....}}\par\vskip 4pt}

\def\ra{\rightarrow}
\newcommand{\blank}[1]{\rule{#1}{0.75pt}}

\newcommand{\prob}[1]{\bigskip\noindent\textbf{#1.} }
\newcommand{\pts}[1]{\,\textsf{#1 pts}}

\newcommand{\probpts}[2]{\prob{#1} \pts{#2} \quad}
\newcommand{\ppartpts}[2]{\textbf{(#1)} \pts{#2} \quad}
\newcommand{\epartpts}[2]{\medskip\noindent \textbf{(#1)} \pts{#2} \quad}


\begin{document}
{\emph{\fontsize{26}{28}\selectfont Math F252\hfill
{\fontsize{32}{36}\selectfont Midterm I}
\hfill Spring 2026}}

\vskip 1.5cm

\emph{\fontsize{18}{22}\selectfont Name:} \quad \underline{\phantom{foo \hspace{9cm} bar}}

\vskip 1cm

{\fontsize{18}{22}\selectfont\emph{Rules:}}

You have {90 minutes} to complete this midterm.  

Partial credit will be awarded, but you must show your work.

%You may have a single handwritten $3 \times 5$ notecard.
Notes, books, calculators, and internet access are not allowed. 

Circle or box your \fbox{FINAL ANSWER} to each question where appropriate.

%If you need extra space, you can use the back sides of the pages.
%Please make it obvious when you have done so.

Turn off anything that might go beep during the exam.

Good luck!

\vfill

\def\emptybox{\hbox to 2em{\vrule height 16pt depth 8pt width 0pt\hfil}}
\def\tline{\noalign{\hrule}}
\centerline{\vbox{\offinterlineskip
{
\bf\sf\fontsize{18pt}{22pt}\selectfont
\hrule
\halign{
\vrule#&\strut\quad\hfil#\hfil\quad&\vrule#&\quad\hfil#\hfil\quad
&\vrule#&\quad\hfil#\hfil\quad&\vrule#\cr
height 3pt&\omit&&\omit&&\omit&\cr
&Problem&&Possible&&Score&\cr\tline
height 3pt&\omit&&\omit&&\omit&\cr
&1&&16&&\emptybox&\cr\tline
&2&&10&&\emptybox&\cr\tline
&3&&10&&\emptybox&\cr\tline
&4&&8&&\emptybox&\cr\tline
&5&&6&&\emptybox&\cr\tline
&6&&28&&\emptybox&\cr\tline
&7&&10&&\emptybox&\cr\tline
&8&&12&&\emptybox&\cr\tline
&Extra Credit&&3&&\emptybox&\cr\tline
&Total&&100&&\emptybox&\cr
}\hrule}}}

\newpage
\prob{1}\ppartpts{a}{2} Sketch the region bounded by the curve $y=1-x^2$ and the $x$-axis.
\vspace{1.8in}

\epartpts{b}{5} Use an integral to calculate the volume of the solid found by rotating the region in part $\textbf{(a)}$ around the \underline{$x$-axis}.
\vfill

\epartpts{c}{5} Use an integral to calculate the volume of the solid obtained by rotating the region in part $\textbf{(a)}$ around the \underline{$y$-axis}.  Use the \emph{shell method}.
\vfill

\epartpts{d}{4} Set up, but \emph{do not evaluate}, an integral for the length of the upper curve of the region sketched in part $\textbf{(a)}$.
\vspace{1in}


\newpage
\probpts{2}{10} A 1-meter long ski pole has linear density \, $\ds \rho(x)=1-\frac{1}{x+3}$ \, grams per centimeter, where one end of the pole is at $x=0$.  Determine the mass of the pole.  Include units.\strut

\vfill

\probpts{3}{10} A 10 centimeter spring requires 20 J of work to stretch the spring to a length of 12 centimeters.  How much work is required to compress the spring to a length of 5 centimeters?  Include units.
\vfill


\newpage
\probpts{4}{8} Evaluate the definite integral, and simplify your answer.

\medskip \noindent  $\ds \int_0^2 x\, e^{2x}\: dx =$
\vfill


\probpts{5}{6}  Set up, \emph{but do not evaluate}, an integral which computes the surface area of the surface of revolution generated by revolving $y=\tan(x^2)$, between $x=0$ to $x=\pi/4$, around the \underline{$x$-axis}.
\vfill


\newpage
\prob{6} Evaluate the indefinite integrals, and simplify your answers.

\epartpts{a}{7} $\ds \int \sin(3x)\sin(4x)\: dx =$
\vfill

\epartpts{b}{7} $\ds \int \sin^2 t \cos^2 t\, dt =$
\vfill


\newpage
\epartpts{c}{7} $\ds \int \ln(5x) \: dx =$
\vfill
\epartpts{d}{7} $\ds \int \frac{2}{x^2 - 3x + 2} \: dx =$
\vfill


\newpage
\probpts{7}{10} Use a trigonometric substitution to evaluate the integral.  Write your final answer in terms of $x$, and simplify it.

\medskip
\noindent $\ds \int \frac{dx}{(4+x^2)^{3/2}} = $
\vfill


\newpage
\prob{8} \ppartpts{a}{4}  For $a$ and $b$ fixed positive numbers, the curve \,$\ds \frac{x^2}{a^2} + \frac{y^2}{b^2}=1$ is an ellipse.  Sketch it.  (\textsl{Hints.}  Find the points along the $x$ and $y$-axes, e.g.~by substituting zero for one of the variables.)

\vspace{2.5in}

\epartpts{b}{8}  Determine the area inside the ellipse by computing an integral.  (\textsl{Hint.}  Start by solving for $y$.)
\vfill


\newpage
\probpts{Extra Credit}{3}  Evaluate, and fully-simplify, the integral you set up in problem \textbf{1 (d)}.
\vfill

\hrulefill

You may find the following \textbf{trigonometric formulas} useful.  Other formulas, not listed here, should be in your memory, or you can derive them from the ones here.
\begin{align*}
\sin(\alpha \pm \beta) &= \sin \alpha \cos \beta \pm \cos \alpha \sin \beta \qquad &
    \sin(ax) \sin(bx) &= \frac{1}{2} \cos((a-b)x) - \frac{1}{2} \cos((a+b)x) \\
\cos(\alpha \pm \beta) &= \cos \alpha \cos \beta \mp \sin \alpha \sin \beta \qquad &
    \sin(ax) \cos(bx) &= \frac{1}{2} \sin((a-b)x) + \frac{1}{2} \sin((a+b)x) \\
 && \cos(ax) \cos(bx) &= \frac{1}{2} \cos((a-b)x) + \frac{1}{2} \cos((a+b)x)
\end{align*}


\clearpage
\newpage

\bigskip
\centerline{\footnotesize \textsc{blank space}}
\vfill

\end{document}